% \chapter{Conclusion and Future work}
\chapter{CONCLUSION AND FUTURE WORK}

\section{Conclusion}

This study investigated the application of Large Language 
Models (LLMs) for grading Thai-language written exam 
responses, focusing on rubric-guided evaluation within the 
framework of Bloom’s Taxonomy. The proposed approach aimed to 
assess the capability of LLMs to provide reliable, 
interpretable, and pedagogically aligned scoring in a 
low-resource linguistic environment.

Despite operating with a modest dataset—comprising 1,620 
graded responses across both rubric and non-rubric conditions 
and a final evaluation on a 55-question test set—the 
experimental outcomes demonstrated promising performance. In 
the rubric-based condition, inter-rater agreement between the 
two expert teachers reached 92.73\%, indicating strong 
consistency when rubric guidance was applied. In contrast, 
non-rubric versus rubric grading by the same teacher yielded 
lower agreement rates (54.55\% for Teacher~1 and 61.82\% for 
Teacher~2), confirming the rubric’s critical role in enhancing 
grading uniformity.

When compared with human grading, the proposed AI-assisted 
system achieved an overall Score Agreement Rate (SAR) of 82.4\% 
and a Cohen’s Kappa ($\kappa$) of 0.78, reflecting 
\textit{substantial agreement} with human raters. Qualitative 
analysis further revealed that 87\% of the AI-generated 
explanations were rated as high quality, demonstrating that 
the LLM not only aligned numerically with human assessments 
but also provided coherent and rubric-consistent reasoning for 
its scoring decisions.

Nevertheless, discrepancies persisted, particularly in 
higher-order cognitive tasks—such as those categorized under 
\textit{Analyzing} and \textit{Creating}—where abstract 
reasoning and creative synthesis posed challenges for both AI 
and human raters. A key source of scoring inconsistency 
stemmed from the rubric’s keyword-based logic. Although the 
rubric was designed to recognize Thai synonyms, the 
morphological richness of the Thai language led to occasional 
mismatches where multiple words conveyed equivalent meanings 
but were not explicitly listed in the rubric. This highlights 
a linguistic constraint that future Thai-language grading 
systems must address through semantic-aware processing or 
embedding-based similarity detection.

The scarcity of Thai written exam data also underscores a 
systemic issue in Thailand’s educational assessment 
landscape—namely, the dominance of multiple-choice testing due 
to the high labor demand of manual grading. This limitation 
has restricted the scale of AI training and evaluation in 
open-ended contexts. In contrast to previous studies conducted 
in high-resource languages such as English or Chinese, this 
research contributes a novel perspective by examining 
rubric-based, LLM-assisted grading in a morphologically 
complex and low-resource language environment. The integration 
of rubric-guided prompting with rule-based algorithmic logic 
resulted in improved consistency, particularly in lower-order 
Bloom’s levels such as \textit{Remembering} and 
\textit{Understanding}.

Bloom’s level proved to be a significant determinant of system 
performance: the AI achieved the highest agreement in recall 
and comprehension-oriented questions, while complex or 
creative responses continued to exhibit variability. Despite 
these challenges, the findings affirm the feasibility of 
deploying AI-assisted grading systems to support Thai 
educators, reduce subjective bias, and enhance the fairness 
and scalability of written assessment.

%%%%%%%%%%%%%%%%%%%%%%%%%%%%%%%%%%%%%%%%%%%%%%%%%%%%%%%%%%%%

\section{Future Work}

Future research will focus on validating the consistency and
accuracy of Bloom’s taxonomy classification using AI-based
methods. In the present study, Bloom levels for all exam
questions were assigned by expert teachers, and these labels
served as the foundation for rubric construction and
evaluation. However, teacher-generated classifications may
contain unintentional human error or subjective variability.

To address this limitation, the next phase of investigation
will examine whether Large Language Models can independently
assign Bloom’s taxonomy levels to Thai written exam questions,
and how closely these AI-derived classifications align with
teacher labels. By comparing teacher–AI agreement patterns,
the system can help identify potential misclassifications,
borderline items, or inconsistencies in cognitive-level
designation.

Additionally, future work will analyze the stability of AI
Bloom-level classification across multiple LLMs to determine
whether certain cognitive levels are more ambiguous or prone to
disagreement. Such analysis will provide insight into which
question types may require clearer wording, more explicit
learning objectives, or rubric refinement.

Finally, AI-assisted Bloom validation may offer a practical
support tool for teachers during assessment design. By
flagging items with low teacher–AI alignment, the system can
help educators verify the appropriateness of question
difficulty and ensure stronger alignment between intended
learning outcomes and actual exam content.

Overall, the next stage of this research will shift toward
using AI not only for grading but also as a mechanism to
evaluate and reinforce the reliability of Bloom’s taxonomy
classification itself, improving transparency and reducing
potential human error in exam development.